% ===========================================================================
\section{Introduction}
\label{sec:intro}
% ===========================================================================

Distributed systems in production environments face a class of faults that are
particularly insidious: \emph{slow faults}~\cite{xinda,gunawi-cbs}. Unlike
crash failures, which are immediately detectable, slow faults manifest as
transient performance degradations---a network link experiencing microsecond
delays, a disk taking milliseconds longer to complete I/O, or a CPU being
temporarily throttled. These degradations are often below monitoring thresholds
yet sufficient to trigger cascading failures in consensus protocols, leader
elections, and replication pipelines~\cite{huang-gray-failure,
  lu-cloud-reliability}.

Industry reports confirm the severity of this problem. Google's Disruptive
Testing Team found that 22\% of production incidents involved degraded (not
failed) components~\cite{google-sre-book}. Microsoft Azure's post-mortem
analysis revealed that gray failures---where a component is in a degraded state
undetected by health checks---account for significant fractions of customer-impacting
outages~\cite{huang-gray-failure}. Datadog's 2023 State of Cloud Monitoring
report documents that median alert thresholds across their customer base are
1--5\,s for network latency anomalies~\cite{datadog-monitoring}, creating a
blind spot for sub-second degradations.

The \xinda{} system~\cite{xinda} (NSDI~'25) made an important contribution by
providing an automated pipeline to systematically explore the impact of slow
faults on distributed systems. However, \xinda{} operates in a
\emph{configured} mode: operators define a severity grid (\eg{},
\texttt{slow-100$\mu$s} through \texttt{slow-1s} in 37 discrete steps), and
the system exhaustively tests each point. This approach has three fundamental
limitations:

\begin{enumerate}[leftmargin=*,topsep=4pt,itemsep=2pt]
  \item \textbf{Grid incompleteness.} Systems whose vulnerability boundaries
        fall outside the pre-defined grid are invisible. Redis Cluster, for
        example, requires 2.7\,s of sustained delay to trigger PFAIL---entirely
        outside \xinda{}'s 100\,$\mu$s--1\,s range. Similarly, filesystem
        fault boundaries for batch-oriented systems (Hadoop, Kafka) often
        exceed standard grid ranges.
  \item \textbf{No minimization.} \xinda{} answers ``does the fault reproduce
        at each grid point?'' (binary yes/no per severity level), but never
        identifies the \emph{minimal} fault configuration---the tightest
        bound on the dangerous region. For systems sensitive at all grid points,
        \xinda{} reports the lowest grid value (100\,$\mu$s) as its bound,
        but the true boundary may be lower still.
  \item \textbf{Grid design burden.} Different systems require different grids
        (Redis needs higher severity values than etcd), defeating the purpose
        of automation and requiring domain expertise from operators.
\end{enumerate}

We present \system{}, a symptom-guided fault reproduction orchestrator that
addresses these limitations. Given a production issue symptom (expressed as a
\emph{symptom oracle}---a declarative specification of log patterns indicating
reproduction), \system{} automatically searches for and \emph{minimizes} a
fault trial that reproduces the symptom. The key insight is that fault
vulnerability boundaries can be efficiently located using binary search over
continuous parameter spaces, rather than exhaustive enumeration over discrete
grids.

\system{}'s core algorithm is a \emph{greedy dimensional minimizer} that
operates in four phases: (1)~eliminate unnecessary faults from multi-fault
trials, (2)~binary-search minimum severity per remaining fault,
(3)~binary-search minimum duration, and (4)~binary-search latest start time.
Each phase uses $O(\log_2 n)$ trials to converge on the boundary in its
dimension, producing a practical upper bound on the minimal fault recipe that
is both human-interpretable and CI-actionable.

We evaluate \system{} across ten widely-deployed distributed systems spanning
five fault types (network, filesystem, CPU, memory, process). Our results show:

\begin{itemize}[leftmargin=*,topsep=4pt,itemsep=2pt]
  \item \system{} uses \textbf{52\% fewer trials} on average than \xinda{}'s
        exhaustive grid across all system--fault configurations, with consistent
        speedups of 1.8--2.4$\times$ in wall-clock time.
  \item \system{} discovers vulnerabilities \textbf{outside fixed grids} in
        7 of 50 configurations---including Redis's 2.7\,s network boundary,
        Hadoop's 1.2\,s filesystem boundary, and Kafka's 850\,ms replication
        boundary.
  \item Consensus-based systems (etcd, ZooKeeper, MongoDB, TiKV, CockroachDB)
        universally exhibit catastrophic failure boundaries at
        \textbf{sub-20\,ms} network delay---250$\times$ below documented
        industry alerting thresholds~\cite{datadog-monitoring,google-sre-book}.
  \item \textbf{Multi-fault interactions} lower individual boundaries by up to
        40\%: combining network + filesystem faults on etcd reduces the network
        boundary from 19.5\,ms to 12.1\,ms.
  \item Results are \textbf{stable across repeated runs}: five independent
        trials per configuration yield coefficient of variation $<$5\% for
        network faults and $<$12\% for filesystem faults.
\end{itemize}

\noindent\textbf{Contributions.} This paper makes three contributions:
\begin{enumerate}[leftmargin=*,topsep=2pt,itemsep=1pt]
  \item The design and implementation of \system{}, a symptom-guided
        orchestrator that produces minimal fault recipes through greedy
        dimensional reduction (\S\ref{sec:design}--\S\ref{sec:impl}).
  \item A comprehensive evaluation across 10 systems $\times$ 5 fault types,
        demonstrating efficiency gains, boundary discovery beyond fixed grids,
        and result stability (\S\ref{sec:eval}).
  \item Empirical characterization of slow-fault danger zones across the
        distributed systems ecosystem, revealing that sub-20\,ms network
        perturbations---routine in cloud environments---universally destabilize
        consensus protocols (\S\ref{sec:findings}).
\end{enumerate}
